% ===================================================================
%  m_W via forward-folding (Matrix Element Method by MC reweighting)
%  WW -> l nu q q-bar  at FCC-ee threshold (ecm 157/160/163)
%  Build:  source /cvmfs/sft.cern.ch/lcg/external/texlive/2024/bin/x86_64-linux  on PATH
%          pdflatex -interaction=nonstopmode mw_forward_fold.tex   (run twice)
% ===================================================================
\documentclass[aspectratio=169,10pt]{beamer}
\usetheme{Madrid}
\usecolortheme{whale}
\setbeamertemplate{navigation symbols}{}
\setbeamertemplate{caption}{\insertcaption}
\setbeamerfont{frametitle}{size=\large}
\setbeamertemplate{frametitle continuation}{}

\usepackage[T1]{fontenc}
\usepackage{lmodern}
\usepackage{amsmath,amssymb,bm}
\usepackage{graphicx}
\usepackage{booktabs}
\usepackage{xcolor}
\usepackage{textcomp}

\graphicspath{{figs/}}
\definecolor{good}{RGB}{0,120,0}
\definecolor{warn}{RGB}{190,0,0}
\newcommand{\mw}{m_W}
\newcommand{\mwhat}{\hat m_W}
\newcommand{\mwhati}{\hat m_{W,i}}
\newcommand{\mqq}{m_{q\bar q}}
\newcommand{\mlv}{m_{\ell\nu}}
\newcommand{\reco}{^{\mathrm{reco}}}
\newcommand{\gen}{^{\mathrm{gen}}}
\newcommand{\GeV}{\,\mathrm{GeV}}
\newcommand{\MeV}{\,\mathrm{MeV}}
\newcommand{\good}[1]{\textcolor{good}{#1}}
\newcommand{\warn}[1]{\textcolor{warn}{#1}}

\title[$\mw$ by forward-folding]{Measuring $\mw$ by forward-folding}
\subtitle{The Matrix Element Method via MC reweighting \\ $WW\to\ell\nu q\bar q$ at the FCC-ee $WW$ threshold}
\author[Defranchis]{M.\ Defranchis}
\institute[CERN]{CERN}
\date{20 June 2026}

\begin{document}

{\setbeamertemplate{footline}{}\begin{frame}\titlepage\end{frame}}

% -------------------------------------------------------------------
\begin{frame}{Outline \& one-sentence summary}
\textbf{Goal.} Measure $\mw$ from semileptonic $WW$ events at the threshold scan
($\sqrt{s}=157/160/163\GeV$, $\sim$92k events / point) by fitting \emph{only} $\mw$.
\vfill
\begin{block}{The whole method in one object}
There is \textbf{one} object: the per-event \emph{folded} likelihood
\[
  L_i(\mw)\;=\;-2\ln P\!\left(\mqq\reco,\;\mlv\reco \,\middle|\, \mw\right)_{\!i},
\]
the resolution-folded double-Breit--Wigner lineshape evaluated on event $i$'s two measured $W$ masses.
\end{block}
\vfill
\textbf{Two read-outs of this one object:}
\begin{itemize}
  \item[(A)] \textbf{Sum} over events $\Rightarrow$ one ensemble $\mwhat$ = \emph{the measurement} (optimal, unbinned ML).
  \item[(B)] \textbf{Per-event} $\mwhati$ (peak of $L_i$) $\Rightarrow$ a data/MC-\emph{checkable} distribution.
\end{itemize}
\vfill
\footnotesize Tableau: motivation $\to$ method (lineshape, why convolution fails, forward-folding)
$\to$ two read-outs $\to$ validation (data/MC, per-event measurement) $\to$ results $\to$ systematics.
\end{frame}

% ===================================================================
\section{Motivation}
% ===================================================================
\begin{frame}{Why a 1-D $\mw$-only fit}
\textbf{The observables} (already in the tree, massless reco objects):
\[
  \mqq\reco=\text{dijet mass},\qquad \mlv\reco=(\ell+\nu)\text{ with }\nu=\vec p_T^{\,\mathrm{miss}}.
\]
\vfill
\textbf{The problem with per-event kinematic fits.} Fitting $\mw$ \emph{plus} per-event nuisances
(jet/lepton scales, ISR, $\nu_z$) opens a near-flat direction:
\begin{itemize}
  \item floating invisible ISR-$z\leftrightarrow\nu_z$ lets the leptonic mass retune to \emph{any} $\mw$;
  \item $\sim$25--60\% of events do not converge (intrinsic, not a solver bug).
\end{itemize}
\vfill
\begin{block}{Fix: stop fitting nuisances per event}
Fit \textbf{only} $\mw$ over the \emph{ensemble}. No per-event free direction $\Rightarrow$
$\sim$100\% of events usable. The detector response is taken \emph{from the MC}, not fit.
\end{block}
\end{frame}

% ===================================================================
\section{The method}
% ===================================================================
\begin{frame}{The lineshape: CC03 double Breit--Wigner $\times$ ISR $\times$ phase space}
\[
  P_{\mathrm{true}}(m_h,m_l;\mw)=\frac{1}{Z(\mw)}\!\int\! d\rho(\sqrt{s'})\;
  \mathrm{BW}(m_h;\mw)\,\mathrm{BW}(m_l;\mw)\,
  \frac{\sqrt{\lambda(s',m_h^2,m_l^2)}}{s'}
\]
\[
  \mathrm{BW}(m;\mw)=\frac{\mw\Gamma_W}{(m^2-\mw^2)^2+(\mw\Gamma_W)^2},\qquad
  \rho(\sqrt{s'})=\text{ISR radiator (Kuraev--Fadin)} .
\]
\vspace{-2mm}
\begin{columns}[T]
\begin{column}{0.6\textwidth}
\includegraphics[width=\linewidth]{lineshape_demo_ecm160}
\end{column}
\begin{column}{0.4\textwidth}\footnotesize
\begin{itemize}
  \item BW peak at $\mw$; the $\sqrt{\lambda}/s'$ \emph{phase space} pulls the observable mean \emph{below} the pole.
  \item $Z(\mw)$ = the kinematic normalisation, computed on-the-fly with Gauss--Legendre ($N=256$) over the
        $\arctan(m^2)$ substitution (the naive GL-24 slopes $\sim\!-0.04/\GeV \Rightarrow +40\MeV$ bias; fixed).
  \item Right panel: the model matches the true gen-mass distribution.
\end{itemize}
\end{column}
\end{columns}
\end{frame}

\begin{frame}{Why the literal (analytic) convolution fails}
\textbf{Plan A} (the obvious one): $P\reco = P_{\mathrm{true}}\otimes K$ with a measured resolution kernel $K$.
\begin{columns}[T]
\begin{column}{0.52\textwidth}
\includegraphics[width=\linewidth]{fit_reco_obs_ecm160}
\end{column}
\begin{column}{0.48\textwidth}\footnotesize
\warn{FAILS closure by $\sim$300\,MeV.} Two unit tests pin the cause:
\begin{enumerate}
  \item \textbf{Mean:} the analytic CC03 marginal mean is $73.8\GeV$ vs the actual gen mean $75.9\GeV$
        ($\Delta=2.1\GeV$): the model shape is wrong.
  \item \textbf{Synth:} even a \emph{stationary-by-construction} kernel fails $\Rightarrow$ convolving a
        mis-shaped $P_{\mathrm{true}}$ with a kernel measured on the real distribution is inconsistent.
\end{enumerate}
The FFT itself is exact (a salvage test convolving the \emph{real} gen histogram closes to $0.23\GeV$).
\textbf{The kernel is also mildly non-stationary} ($\langle\Delta\mqq\rangle$ drifts with true mass).
\end{column}
\end{columns}
\vfill
\centering\textbf{$\Rightarrow$ Don't model the folding analytically. Take it from the MC.}
\end{frame}

\begin{frame}{The fix: forward-folding = MEM by MC reweighting}
\textbf{Idea.} The MC events already carry the \emph{real} per-event gen$\to$reco smearing.
Build the reco template at a trial $\mw$ by \emph{reweighting} the actual MC reco masses:
\[
  w_i(\mw)=\frac{P_{\mathrm{true}}(m\gen_i;\mw)}{P_{\mathrm{true}}(m\gen_i;m_{W0})},\qquad
  T(m\reco;\mw)=\sum_i w_i(\mw)\,\big[\,m\reco_i\in\text{bin}\,\big].
\]
\vspace{-1mm}
\includegraphics[width=\linewidth]{reweight_demo_ecm160}
\vspace{-1mm}
\footnotesize
(A) the weight tilts the \emph{fixed} sample; (B) the gen peak slides to the trial $\mw$;
(C) the reco template slides $\Rightarrow$ the fit picks the $\mw$ whose template matches the data.
\textbf{$k$-fold} (build template from other folds) removes self-use $\Rightarrow$ all $N$ events measured, full-sample $\sigma$.
\end{frame}

\begin{frame}{This is literally the Matrix Element Method}
\[
  P(m\reco|\mw)=\frac{1}{\sigma(\mw)}\!\int \underbrace{|M(\mw)|^2}_{\mathrm{BW\cdot BW\cdot PS}}\;
  \underbrace{f(\sqrt{s'})}_{\mathrm{ISR}}\;\underbrace{W(m\reco|\mathrm{true})}_{\mathrm{transfer\ fn.}}\;d(\mathrm{true})
\]
\begin{itemize}
  \item The ``convolution'' was MEM with an \emph{analytic} transfer function $W$ (the kernel) -- it failed.
  \item Forward-folding uses the \textbf{MC events as samples of $W(m\reco|\mathrm{true})$}, reweighted by
        $|M(\mw)|^2/|M(m_{W0})|^2$ = the standard MEM-by-reweighting / template-morphing variant.
\end{itemize}
\vfill
\begin{block}{Two approximations (both calibratable)}
(1) project the full parton phase space onto the two masses;\quad
(2) CC03 double-BW instead of the full 4-fermion ME $\Rightarrow$ a $\sqrt{s}$-dependent lineshape bias
($-93/-120/+27\MeV$), removed by full-4f reweighting.
\end{block}
\end{frame}

\begin{frame}{The 2-D resolution kernel (diagnostic input)}
\centering
\includegraphics[height=0.78\textheight]{kernel_diag_ecm160}\\
\footnotesize The measured 2-D $(m\reco-m\gen)$ response and its (mild) non-stationarity. Used for diagnostics
\& the synthetic unit test; the \emph{measurement} uses the per-event MC smearing directly (no kernel assumption).
\end{frame}

% ===================================================================
\section{Two read-outs}
% ===================================================================
\begin{frame}{Read-out (A): sum the curves = the measurement}
\centering
\includegraphics[width=\linewidth]{sum_of_curves_ecm160}
\vfill
\footnotesize
Each event's $L_i(\mw)$ is \emph{broad} (left). Their \emph{sum} is razor-sharp (right):
$\mwhat=\arg\min_{\mw}\sum_i L_i(\mw)$, $\sigma_{\mw}\!\approx\!7\MeV$. Informative events are auto-weighted by curvature.
\textbf{This is the optimal, unbinned maximum-likelihood number.}
\end{frame}

\begin{frame}{Read-out (B): the per-event $\mwhati$ distribution}
\begin{columns}[T]
\begin{column}{0.62\textwidth}
\includegraphics[width=\linewidth]{perevent_mW_ecm160}
\end{column}
\begin{column}{0.38\textwidth}\footnotesize
$\mwhati=\arg\max_{\mw}P_{\mathrm{folded}}(\text{masses}_i|\mw)$ -- a \emph{fixed} per-event observable
(no nuisance fit $\Rightarrow$ no flat direction, defined for $\sim$100\% of events).
\begin{itemize}
  \item (A) its \emph{distribution} is a data/MC-checkable observable (validates the folding);
  \item (B) its weighted mean \emph{shifts with the true} $\mw$ (slope $\sim$0.75) $\Rightarrow$ calibratable.
\end{itemize}
This is the top-mass ``ideogram'' / per-event-likelihood method.
\end{column}
\end{columns}
\end{frame}

% ===================================================================
\section{Validation}
% ===================================================================
\begin{frame}{Staged validation -- each stage isolates one piece}
\centering\footnotesize
\begin{tabular}{lll}
\toprule
Stage & What it tests & Result @160 \\
\midrule
\texttt{gen\_pe}   & lineshape + normalisation (per-event exact $\sqrt{s'}$) & 80.205, 100\% events, pull 1.06 \\
\texttt{gen\_marg} & marginalisation over the ISR radiator & 80.259 (agrees with \texttt{gen\_pe} family) \\
\texttt{reco} (conv.) & analytic convolution & \warn{79.9 -- FAILS by $\sim$300\,MeV} \\
\texttt{fold}      & forward-folding (the method) & \good{80.235, $\sim$100\% events, pull 0.95} \\
\bottomrule
\end{tabular}
\vfill
\normalsize
\begin{block}{Estimator gotcha (caught by an identity unit test)}
Gaussian-smoothing a sharp template biases the morphing MLE by $-170\MeV$. The bias-free operating point
is \textbf{fine bins (0.5\,GeV) + light/no smoothing}. The identity test (\texttt{FOLD\_TGT=gen}, which must
reproduce $m_{W0}$) catches it; the reco fold sits \emph{on} this estimator floor.
\end{block}
\end{frame}

\begin{frame}{Data/MC validation -- pre/post-fit projection + pull (ecm160)}
\centering
\includegraphics[width=0.72\linewidth]{validation_prepostfit_ecm160}\\[0.5mm]
{\footnotesize
Pseudo-data (40\% held-out) vs an MC template from the \emph{disjoint} 60\%. \textbf{Post-fit} describes the data
($\chi^2/\mathrm{ndf}\!\approx\!1$: 51/54, 61/54); a \textbf{mis-set} $\mw$ ($+0.5\GeV$) gives a coherent pull at
the $W$ peak ($\chi^2=344/552$) $\Rightarrow$ the pull panel \emph{flags} a wrong $\mw$ / a data--MC disagreement.\par}
\end{frame}

\begin{frame}{Data/MC validation -- ecm157 \& ecm163}
{\footnotesize
Post-fit $\chi^2/\mathrm{ndf}$ ($\mqq$,$\mlv$): \;157: 56/54,\,65/54;\;\; 160: 51/54,\,61/54;\;\; 163: 51/54,\,56/54
-- all $\approx1$. Mis-set $\mw$ gives $\chi^2=200$--$552$ $\Rightarrow$ clear discriminating power everywhere; the
\emph{same} overlay a template fit calls its ``check'' (unbinned ML separates it into this plot -- nothing lost).\par}
\vfill
\centering
\includegraphics[width=0.49\linewidth]{validation_prepostfit_ecm157}\hfill
\includegraphics[width=0.49\linewidth]{validation_prepostfit_ecm163}\\[0.5mm]
{\scriptsize ecm157 (left), ecm163 (right): pseudo-data (40\% held-out) vs the MC template from the disjoint 60\%.}
\end{frame}

\begin{frame}{The per-event read-out as an actual measurement}
\centering
\includegraphics[width=\linewidth]{perevent_measurement_ecm160}
\vfill
\footnotesize
\textbf{(A)} the per-event $\mw$ likelihood density = the \emph{stacked per-event posterior} (each event's smooth
likelihood curve summed $\Rightarrow$ a smooth ideogram); reweighting to true $\mw=\mu$ \emph{shifts} it
(dotted/dashed). \textbf{(B) Closure:} slope $1.000$, residual $<\!1\MeV$ over $\pm0.4\GeV$.
\textbf{(C) Toys:} bias $\le1\MeV$, $\sigma_{\mw}=8.1\MeV$ vs ensemble $7.1\MeV$ $\Rightarrow$ the transparent,
histogram-native read-out costs only $\sim$15--20\%.
\end{frame}

% ===================================================================
\section{Results}
% ===================================================================
\begin{frame}{Headline: closure \& precision (forward-fold, $k$-fold)}
\begin{center}\includegraphics[width=0.9\linewidth]{fold_final_results}\end{center}
\vspace{1pt}
{\footnotesize
\textbf{Left:} $\mwhat=80.289/80.235/80.385$; closure vs gen target \good{$+3/-24/-21\MeV$}; fold sits on the
(calibratable) estimator floor. \textbf{Right:} $\sigma_{\mw}=7.6/7.1/10.1\MeV$; \textbf{2-D $\sim$40\% tighter
than $\mqq$-only}, at the independent-combination floor (score corr $\approx0$).\par\smallskip
\textcolor{warn}{Not a test of CC03:} the model bias \emph{cancels} in closure (same model both sides) $\Rightarrow$
residual = estimator floor, not model failure. The CC03-vs-4f inadequacy is the \emph{separate absolute} offset
(gen target vs PDG $=-93/-120/+27\MeV$); fold tolerates the shape error the convolution could not.\par}
\end{frame}

\begin{frame}{The decision: which read-out? -- BOTH}
\centering
\includegraphics[width=0.96\linewidth]{readout_compare}
\vfill
\footnotesize
They are \emph{two read-outs of the same likelihood}, not competing methods.
\textbf{Ensemble} = the headline measurement (optimal, $\sim$100\% events, $\sigma\,7$--$10\MeV$).
\textbf{Per-event template fit} = the validation / control-region read-out (transparent, data/MC-checkable,
$+$15--20\% on $\sigma$; closure slope 1.00, bias $<$ few MeV across $\pm0.4\GeV$).
\end{frame}

\begin{frame}{Results summary table}
\centering
\begin{tabular}{lccc}
\toprule
 & ecm157 & ecm160 & ecm163 \\
\midrule
events (usable)                 & 92.5k & 92.2k & 92.2k \\
\addlinespace
ensemble $\mwhat$ [GeV]         & 80.289 & 80.235 & 80.385 \\
ensemble $\sigma_{\mw}$ [MeV]   & 7.6 & 7.1 & 10.1 \\
closure vs gen target [MeV]     & \good{$+3$} & \good{$-24$} & \good{$-21$} \\
\addlinespace
per-event fit: closure slope    & 1.000 & 1.000 & 1.000 \\
per-event fit: bias [MeV]       & $+0$ & $-1$ & $-1$ \\
per-event fit: $\sigma_{\mw}$ [MeV] & 8.5 & 8.6 & 11.7 \\
per-event / ensemble $\sigma$   & 1.12$\times$ & 1.21$\times$ & 1.16$\times$ \\
\addlinespace
post-fit $\chi^2/$ndf ($\mqq$,$\mlv$) & 56/54,\,65/54 & 51/54,\,61/54 & 51/54,\,56/54 \\
\bottomrule
\end{tabular}
\vfill
\footnotesize $\mw$ is the \emph{only} free parameter; pull $\approx1.0$; seed- and machine-robust.
\end{frame}

\begin{frame}{Statistical reach: the unbinned ML extrapolates like an Asimov fit}
\begin{columns}[c]
\begin{column}{0.56\textwidth}
\includegraphics[width=\linewidth]{asimov_scaling}
\end{column}
\begin{column}{0.44\textwidth}\footnotesize
\textbf{Q: can we project to a larger sample, as a template Asimov fit allows?} \textbf{Yes -- identically.}
\vfill
The ensemble test statistic is a \emph{sum over events},
\[ -2\Delta\ln L(\mw)=\textstyle\sum_i L_i(\mw), \]
so the Fisher information is additive, $I(\mw)=\sum_i I_i\propto N$, and $\sigma_{\mw}=1/\sqrt{I}\propto 1/\sqrt{N}$.
\vfill
\begin{itemize}
  \item Two equivalent extrapolations: (i) rescale $\sigma\!\to\!\sigma\sqrt{N/N'}$; (ii) the \emph{Asimov}
        curve = $N'\!\times$ the per-event expected $-2\ln L$ -- read its curvature, no toys.
  \item Verified here: subsampling the MC follows $1/\sqrt{N}$, and the curvature $\sigma$ \emph{equals} the
        bootstrap $\sigma$ (asymptotic regime).
  \item $\Rightarrow$ statistical-only reach $\sigma\!\approx\!2.2/0.7/0.2\MeV$ at $10^{6}/10^{7}/10^{8}$ events.
        ``Unbinned'' does \emph{not} forfeit Asimov projection -- the per-event likelihood already \emph{is} the
        Asimov ingredient.
\end{itemize}
\end{column}
\end{columns}
\end{frame}

% ===================================================================
\section{Systematics \& outlook}
% ===================================================================
\begin{frame}{Anatomy of the closure: it is the estimator floor, not modeling}
\centering
Decompose the matched closure:
$\;\text{closure}=\underbrace{\text{floor}}_{\text{identity fold}-\text{gen target}}+\underbrace{\text{smear}}_{\text{fold(reco)}-\text{identity fold}}$.
\\[0.5mm]
\includegraphics[width=0.80\linewidth]{closure_binning_floor}\\[0.5mm]
{\scriptsize
\textbf{Binning probe} (\texttt{FOLD\_BIN} $0.25\!\to\!1.0$): the \warn{floor swings $\sim$100--145\,MeV} (left)
while \good{$\mathtt{gen\_pe}$/$\mathtt{gen\_marg}$/gap are bit-invariant} (right). Variance share of the closure
scatter: \textbf{floor $0.71/0.63$} vs smear $0.14/0.04$ $\Rightarrow$ the ECM-dependent closure \emph{scatter} is a
\good{calibratable binned-template floor}, \textbf{not} a radiator/threshold/4f \emph{modeling} effect (revising the
earlier hypothesis). Fine bins shrink it; genuine smearing residual $\sim\!3$--$8\MeV$, binning-stable.\par}
\end{frame}

\begin{frame}{The grid$-$exact gap = a threshold-driven marginalisation bias}
\begin{columns}[c]
\begin{column}{0.60\textwidth}\centering
\includegraphics[width=\linewidth]{closure_anatomy}
\end{column}
\begin{column}{0.40\textwidth}\footnotesize
The real modeling effect is the gen-target gap
$\mathtt{gen\_marg}-\mathtt{gen\_pe}=+70/+54/+11\MeV$.
\smallskip

\textbf{(B) Width ladder:} scaling the $\sqrt{s'}$ spread $f\!:\!1\!\to\!0$ closes the gap; its \emph{spread} term
\warn{flips sign across $2\mw=160.76$} ($+41/+38\!\to\!-27$) $\Rightarrow$ phase-space curvature / threshold.
\smallskip

\textbf{(C)} grows with off-shellness $\delta=\sqrt{s'}-(\mqq+\mlv)$; \textbf{(D)} level set by the $\sqrt{s'}$ regime.
\smallskip

\textbf{Decision:} \texttt{FOLD\_W=exact} (per-event $\sqrt{s'}$) removes the gap by construction $\Rightarrow$
principled for the \emph{absolute} $\mw$ (fine bins minimise its floor); for closure-calibrated \emph{relative} $\mw$
it cancels.
\end{column}
\end{columns}
\end{frame}

\begin{frame}{Systematics \& caveats}
\begin{itemize}
  \item \textbf{Reweighting weight (grid vs exact):} \good{resolved} -- grid carries a $+11$ to $+70\MeV$
        \emph{radiator-marginalisation} bias (threshold-driven, spread term flips at $2\mw$); use
        \texttt{FOLD\_W=exact} for the absolute $\mw$ (prev.\ slides).
  \item \textbf{Estimator floor $=$ the binning systematic:} the identity-test residual swings $\sim$100--145\,MeV
        coarse$\to$fine and is the dominant piece of the closure scatter; \good{calibratable}, minimised by fine bins
        ($\le0.25$); quote its bin/sm/seed stability. \warn{Avoid} the coarse-bin $+$ heavy-smooth corner.
  \item \textbf{CC03 vs full 4f:} the \emph{separate, larger} $\sqrt{s}$-dependent absolute bias
        ($-93/-120/+27\MeV$); removed by full-4f reweighting.
  \item \textbf{Data/MC resolution (the leading exp.\ systematic):} the validation plot \emph{is} the handle --
        distort the MC resolution and re-fit to propagate it.
\end{itemize}
\vfill
\footnotesize All absolute-$\mw$ items need an MC/4f calibration; the \emph{closure} and \emph{precision} above are robust.
\end{frame}

\begin{frame}{Summary \& next steps}
\begin{block}{Summary}
\begin{itemize}
  \item One folded likelihood $L_i(\mw)$, two read-outs: ensemble (optimal) + per-event distribution (checkable).
  \item Forward-folding (MEM by MC reweighting) \good{works}: closure $+3/-24/-21\MeV$, $\sigma_{\mw}=7$--$10\MeV$,
        $\sim$100\% events, pull $\approx1$. The analytic convolution fails; the MC transfer function fixes it.
  \item Full data/MC validation (pre/post-fit + pull) and a closing per-event template-fit measurement (slope 1.00).
\end{itemize}
\end{block}
\begin{block}{Next}
\begin{enumerate}
  \item \good{Done:} the ECM-dependent closure scatter $=$ a calibratable binning floor (\emph{not} modeling); the
        grid$-$exact gap $=$ a threshold marginalisation bias $\Rightarrow$ \texttt{FOLD\_W=exact}.
  \item Reduce the exact estimator floor (KDE / unbinned template / ISR-tail weight regularisation) $\to$ quote the
        exact absolute headline; then full-4f reweighting to kill the (larger) CC03 bias.
  \item Data/MC resolution systematic via the validation plot; C++ port (lineshape $+\log Z$ exist; needs the
        gen-reweight $+$ reco-template stage).
\end{enumerate}
\end{block}
\end{frame}

{\setbeamertemplate{footline}{}\begin{frame}[plain]\centering\Large Backup\end{frame}}

\begin{frame}{Backup: ensemble fit + bootstrap (ecm160)}
\centering
\includegraphics[width=0.92\linewidth]{fit_fold_obs_ecm160}\\
\footnotesize Left: the ensemble $-2\Delta\ln L(\mw)$ parabola. Right: the bootstrap $\mwhat$ distribution
(stat.\ $\sigma$, pull cross-check).
\end{frame}

\begin{frame}{Backup: per-event $\mwhati$ distributions, 3 energies}
\centering
\includegraphics[width=0.96\linewidth]{fold_mW_distributions}\\
\footnotesize The per-event $\mwhati$ distribution at each energy -- the data/MC-checkable companion observable.
\end{frame}

\end{document}
